\section*{Your turn}

Everything in this chapter was measured against two services you can start in
two terminals. Check out tag \code{ch08} of the companion repository and work
from its root. You need the .NET SDK, Docker, and \code{npm ci} to have been run
once. Ports 5100 and 5101 have to be free.

\begin{minted}{text}
docker compose up -d mosaic-db
dotnet run --project src/Mosaic.Catalog   # :5101
dotnet run --project src/Mosaic.Api       # :5100
\end{minted}

The first start of each creates and seeds its own database and says so.
Catalog logs \code{Seeded 25 products into a newly created database.} and Mosaic
logs a line naming prices, reviews, customers and orders. If either sits there
without logging a seed line, it is waiting on PostgreSQL; \code{docker compose
ps} should show \code{mosaic-db} healthy.

\begin{enumerate}
  \item \textbf{Read both contracts.} Send \code{\{ \_service \{ sdl \} \}} to
  port 5101 and then to port 5100, and put the two \code{Product} blocks side by
  side. Nine distinct fields between them, and the only one declared twice is
  the key. Then find the two differences that are not fields: one type
  implements \code{Node} and the other does not, and only one of them has
  \code{ProductCategory} in its schema at all.

  \item \textbf{Assemble one product by hand.} Ask Catalog for
  \code{\{ products \{ id title \} \}} and copy the first \code{id}. Send it to
  Mosaic as a representation, asking for \code{price} and \code{reviews(first:
  2)}. You have just executed a query plan with your hands. Write down how many
  requests it took and which service you had to ask first, because that ordering
  is not a convention, it is the only order that works.

  \item \textbf{Break the key three ways.} Send Catalog a representation whose
  \code{id} is the raw \code{Guid} \code{a0000000-0000-4000-8000-000000000001},
  then one with the string \code{not-a-key}, then a valid product key with
  \code{\_\_typename} changed to \code{Customer}. The first two answer null. The
  third does not answer null, and working out why from the response alone is the
  exercise.

  \item \textbf{Watch the batching, or fail to.} Turn on Entity Framework Core's
  command logging for Catalog:

  \begin{minted}[fontsize=\footnotesize]{text}
$env:ASPNETCORE_ENVIRONMENT = 'Development'
${env:Logging__LogLevel__Microsoft.EntityFrameworkCore.Database.Command} = `
    'Information'
  \end{minted}

  and restart it. The brace syntax is not optional: the category name contains
  dots and PowerShell needs \code{\$\{env:...\}} to set it. Now send all
  twenty-five product keys in one \code{\_entities} call and count the
  \code{Executed DbCommand} lines. Then send them one call at a time and count
  again. The answers are identical and the counts are not.

  \item \textbf{Put the reference resolver in the wrong place.} In
  \code{Mosaic.Catalog}, move the attributed method off the \code{Product}
  record and onto \code{ProductNode}, the type extension class beside it. It
  compiles. Restart, ask for \code{\_service}, and find the field that
  appeared. Then send any representation and read what comes back. Put it back
  afterwards.

  \item \textbf{Take the type name check out and find that nothing happens.} In
  \code{ProductKey}, return the parsed value without comparing
  \code{nodeId.TypeName}. Restart Catalog and send it a representation whose
  \code{id} is a \emph{customer} key, taken from any review's author. You get a
  null, which is what you got before. Now work out what would have to be true
  about Mosaic's identifiers for the two answers to differ, and decide whether
  you would ship the check anyway. My answer is in
  section~\ref{sec:ch08-the-key-you-already-designed}; disagree with it if you
  can say why.

  \item \textbf{Ask Mosaic for something that left.} Send
  \code{\{ products \{ id \} \}} to port 5100 and read the error carefully. Then
  ask yourself what would have told you, before you deployed, that the field had
  gone. The answer is chapter~\ref{ch:ci-cd-for-a-federated-graph}, and it is
  worth having the question first.

  \item \textbf{Compose them, and then break it.} From the repository root:

  \begin{minted}[fontsize=\footnotesize]{text}
npx wgc router compose -i federation/mosaic.yaml -o supergraph.json
  \end{minted}

  It succeeds. Now open either \code{Program.cs} and drop the
  \code{registerNodeInterface} argument, so that the call is a plain
  \code{AddGlobalObjectIdentification()} again. Re-export that service's schema
  over its committed file:

  \begin{minted}[fontsize=\footnotesize]{text}
dotnet run --project src/Mosaic.Catalog -- \
    schema export --output ../../schema/catalog.graphql
  \end{minted}

  The doubled-back path is not a typo. \code{schema export} resolves
  \code{--output} against the \emph{project} directory rather than the one you
  are standing in, so the obvious spelling of that command writes a file inside
  \code{src/Mosaic.Catalog/} and leaves the committed schema untouched. Compose
  again and read the error. Then put it back and do the same for the two cost
  options, and again for the shareable page cursor class. Three settings, three
  errors, and no test of a single service would have found any of them.

  \item \textbf{Run the gate.} \code{pwsh scripts/verify.ps1}. It starts both
  services, checks both published schemas against the committed files, asks the
  chapter's question in two halves, and composes the pair. Then change
  \code{\$ExpectedSqlCommandCount} from 2 to 3 and run it again, to see what a
  stale number in this chapter would look like to the build.
\end{enumerate}

If you do only one of these, do exercise~8. Everything else here can be worked
out from one service at a time, and that is the last chapter in this book where
that is true.
